% Generated draft: requires mathematical and source review.
\documentclass[11pt]{article}
\usepackage[margin=25mm]{geometry}
\usepackage{fontspec}
\setmainfont{Latin Modern Roman}
\usepackage{amsmath,amssymb,mathtools,unicode-math}
\setmathfont{Latin Modern Math}
\usepackage{enumitem,needspace,xurl,hyperref}
\setlength{\parindent}{0pt}
\setlength{\parskip}{5pt}
\setlength{\emergencystretch}{3em}
\newcommand{\N}{\mathbb{N}}
\newcommand{\Z}{\mathbb{Z}}
\newcommand{\Q}{\mathbb{Q}}
\newcommand{\R}{\mathbb{R}}
\newcommand{\C}{\mathbb{C}}
\newcommand{\F}{\mathbb{F}}
\newcommand{\A}{\mathbb{A}}
\newcommand{\PP}{\mathbb{P}}
\newcommand{\E}{\mathbb{E}}
\newcommand{\eps}{\varepsilon}
\newcommand{\dd}{\,\mathrm{d}}
\begin{document}
\section*{UnsolvedMath ID: 30005440}
\textbf{Source problem code:} catalogue entry
\subsection{Definitions and mathematical statement}
Let \(p\) be a prime. A finite \(p\)-group is a group of order \(p^n\) for some integer \(n\geq 0\). Under the conventional interpretation of a brace as a left brace with abelian additive group, a brace is a set \(B\) with operations \(+\) and \(\circ\) such that \((B,+)\) is an abelian group, \((B,\circ)\) is a group, and
\[
 a\circ(b+c)=a\circ b-a+a\circ c
 \qquad(a,b,c\in B).
\]
Its multiplicative group is \((B,\circ)\). A skew left brace satisfies the same identity but need not have abelian additive group; thus the distinction between braces and skew braces is material here.

For a finite \(p\)-group \(G\), define the realization property
\[
\mathcal R(G)\quad\Longleftrightarrow\quad
\text{there exists a finite left brace }(B,+,\circ)
\text{ with }(B,\circ)\cong G.
\]
Equivalently, one asks whether an abelian group operation \(+\) can be placed on the underlying set of \(G\) so that its given multiplication and \(+\) satisfy the displayed brace identity.

\textbf{Statement under this interpretation.} Determine which finite \(p\)-groups satisfy \(\mathcal R(G)\), as \(p\) ranges over the primes. Give a characterization, or necessary or sufficient group-theoretic conditions, for this realization property. In particular, determine a smallest finite \(p\)-group failing \(\mathcal R\), with the following qualification about the intended minimum. If size means group order and \(p\) is fixed, the numerical target is
\[
 m_p=\min\bigl\{|G|:G\text{ is a finite }p\text{-group and }\neg\mathcal R(G)\bigr\},
\]
whenever this set is nonempty, together with an example attaining the minimum. If the minimum is intended across all primes, the corresponding target is
\[
 m=\min\bigl\{|G|:G\text{ is a finite }p\text{-group for some prime }p,
 \ \neg\mathcal R(G)\bigr\},
\]
again conditional on nonemptiness. The supplied record does not specify which minimum is intended, and these alternatives must not be silently identified.

\textbf{Unresolved source qualifications.} The source excerpt asks for a ``minimal counterexample'' without supplying its preceding conjectural referent. The interpretation above takes a counterexample to mean a finite \(p\)-group that cannot be realized as the multiplicative group of a brace, following the supplied clean statement. The excerpt does not itself specify the left/right convention for braces or explicitly state that the additive group is abelian; the definition above is therefore a stated interpretation requiring source review.

The excerpt also reports that a skew brace factorized into two right nilpotent skew braces need not be right nilpotent, and that the question remains open when both factors are ideals of the larger brace. This is retained as supplied contextual material. The excerpt does not define its factorization convention or its right-nilpotence convention, and it does not clearly establish how this contextual question relates to the multiplicative-group realization problem. It cannot be promoted to a precise additional assertion without that information.
\subsection{Short English statement}
Determine which finite \(p\)-groups are multiplicative groups of braces, and seek a characterization, necessary or sufficient conditions, and a smallest nonrealizable example. The draft explicitly uses left braces with abelian additive group. The intended brace convention, the referent of ``counterexample,'' and whether the minimum is for a fixed prime or across all primes require review. The accompanying skew-brace factorization question is preserved as unresolved context.
\subsection{Sources}
\begin{itemize}
\item Local UnsolvedMath registry: permanent ID \(30005440\), title \emph{Multiplicative p-Groups of Braces}, source problem code \texttt{catalogue entry}, in the supplied Oberwolfach Reports Open Problems collection. No individual proposer is identified; the supplied proposed year is 2023.
\item Supplied primary reference: \emph{Mini-Workshop: Skew Braces and the Yang--Baxter Equation} (2023), Oberwolfach Reports, pp.~537--563, DOI \texttt{10.4171/owr/2023/9}. The extracted passage includes the page marker 549. \url{https://doi.org/10.4171/owr/2023/9}. This is also the supplied literature-source link labeled ``original OWR report.''
\item Supplied literature source: Bachiller, \emph{Counterexample to a conjecture about braces}, DOI \texttt{10.1016/j.jalgebra.2016.01.011}. \url{https://doi.org/10.1016/j.jalgebra.2016.01.011}.
\item Literature assessment dated \texttt{2026-08-22}: the supplied curation labels the problem open and reports positive results through order \(p^4\), known counterexamples, and no requested characterization or minimum. These are supplied dated curation claims, not independently verified mathematical or literature conclusions. The background explicitly requests independent verification of time-sensitive status and cited claims.
\item The registry labels the statement \texttt{corrected\_verified} and says that the official report was checked directly and the clean statement restores definitions and removes extraction artifacts. These are supplied verification claims; no independent source visit or verification is asserted in this draft.
\end{itemize}
\end{document}
